% Targeted second reading, 2026-11-18, GPT-6 Astra Pro.
% Pro mode; a separate effort setting was not exposed.
% This is an unnumbered proof expansion for B.B8.3, not a new priority claim.
\subsection{Second reading of the norm-polynomial reconstruction}

The finite-grid and equivariance arguments in B.B8.3 are valid under
the arithmetic inputs B.B8.1--B.B8.2 and the pinned-point-map input
quoted there. The following expands the global continuation step.
It supplies the details needed when holomorphic and antiholomorphic
root branches meet. No continuity of the original point map is used.

\paragraph{Finite-grid recovery.}
For a monic polynomial $P(Z)=\sum_{j=0}^d c_jZ^j$, $c_d=1$, write
\[
 B(Z,W)=P(Z)\sum_{j=0}^d\overline{c_j}W^j.
\]
After $Z=r+is$, $W=r-is$, the polynomial $B$ has degree at most
$2d$ in each of $r,s$. Its $(2d+1)^2$ values on the stated grid
therefore determine it by successive univariate interpolation.
The coefficient of $Z^kW^d$ is $c_k$, without a phase ambiguity,
because $c_d=1$. Hence the reconstructed coefficients are fixed
algebraic linear functionals of the squared moduli. Translation
and multiplication of the grid by a nonzero algebraic scalar are
invertible algebraic changes of the interpolation coordinates.
In particular the grid can avoid any prescribed finite set.

\paragraph{The differential identity.}
On the common complex Zariski open set $\Omega$, interpolation
produces single-valued real-analytic coefficients agreeing with the
actual fibre polynomial at every algebraic parameter. For each
algebraic $a$, the norm formula on its own open set consequently
holds as a real-analytic identity, since algebraic parameters are
Euclidean dense. On $P(t,a)\ne0$, pluriharmonicity gives
\[
 Q_{jk}(a)+\overline{Q_{kj}(a)}=0,\qquad
 Q_{jk}(a)=\frac{P_{t_j\bar t_k}}{P}
          -\frac{P_{t_j}P_{\bar t_k}}{P^2}.
\]
For a fixed algebraic $a$, this identity extends through the
auxiliary hypersurfaces wherever $P(t,a)\ne0$, by continuity of
its real-analytic left side. At a fixed $t$, it holds for every
algebraic $a$ outside the finite root set, and hence for every
complex $a$ outside that set. The rational holomorphic function
$Q_{jk}$ is also antiholomorphic there, so it is constant. It tends
to zero at infinity because the leading coefficient of $P$ is the
constant one. Thus $Q_{jk}=0$ for all $j,k$.

\paragraph{Local algebraic separation of the root types.}
At a generic real-analytic germ, complexify $\bar t$ to independent
variables $w$. The coefficients of $P$ belong to a finite algebraic
extension of $\overline{\mathbf Q}(t,w)$. In a finite splitting
field the derivations $\partial_{t_j}$ and $\partial_{w_k}$ extend
uniquely, since the characteristic is zero. Decompose $P$ into its
distinct roots $\lambda_\nu$ with positive multiplicities $m_\nu$.
The double and simple partial-fraction poles of
\[
 \partial_{t_j}\partial_{w_k}\log P(t,w,a)=0
\]
give respectively
\[
 m_\nu(\partial_{t_j}\lambda_\nu)
        (\partial_{w_k}\lambda_\nu)=0,
 \qquad
 m_\nu\partial_{t_j}\partial_{w_k}\lambda_\nu=0.
\]
If one $t$-derivative of a root is nonzero, every $w$-derivative
is zero; if one $w$-derivative is nonzero, every $t$-derivative
is zero. This is an identity in a field, so no point-dependent
choice of type is involved. A root with both sets zero is constant.

A root killed by all $w$-derivations is algebraic over
$\overline{\mathbf Q}(t)$. Here is the elementary differential
algebra argument. Take its monic minimal polynomial over
$\overline{\mathbf Q}(t,w)$. Derivations commute with the
automorphisms of a separable normal closure, so all its conjugates
are also killed. Its coefficients, the elementary symmetric
functions in those conjugates, belong to the simultaneous constant
field of the $w$-derivations, namely $\overline{\mathbf Q}(t)$.
The same argument applies with $t$ and $w$ exchanged. A root killed
by both sets is algebraic over $\overline{\mathbf Q}$ and hence
belongs to that algebraically closed field.

We have therefore obtained monic algebraic coefficient germs
\[
 P(t,Z)=H_0(t,Z)A_0(\bar t,Z),
\]
where the roots of $H_0$ are holomorphic or constant, and those of
$A_0$ are antiholomorphic and nonconstant. Repeated roots stay with
their full multiplicities. A holomorphic nonconstant algebraic
germ cannot equal an antiholomorphic one: complexifying such an
equality makes all derivatives in both independent variable sets
zero. This observation gives uniqueness of the division into the
two root types, with constants assigned to $H_0$.

\paragraph{Global continuation, including collisions of opposite types.}
Remove from $\Omega$ a complex algebraic hypersurface containing
the poles and branch loci of the finite algebraic coefficient
extensions of $H_0(t,Z)$ and of the conjugate of
$A_0(\bar t,Z)$. Call the resulting connected complex Zariski
open set $\Omega'$. Each coefficient germ can now be continued
along every path in $\Omega'$. Their product continues to $P$,
by uniqueness of continuation of real-analytic germs. This does
not require the roots of $H_0$ to avoid the roots of $A_0$ on the
real diagonal; continuation of the algebraic coefficients of
the two factors is defined independently of those collisions.

After continuation around any loop, let the resulting monic
factors be $H_1(t,Z)$ and $A_1(\bar t,Z)$. They satisfy
\[
 H_1(t,Z)A_1(\bar t,Z)=H_0(t,Z)A_0(\bar t,Z).
\]
Continuation preserves the vanishing of the relevant derivations
and cannot turn a nonconstant algebraic germ into a constant one.
Uniqueness of the two root types, just proved, forces $H_1=H_0$
and $A_1=A_0$. Thus every coefficient of each factor is single
valued on $\Omega'$; the first is holomorphic there, and the
second becomes holomorphic after conjugating the coordinates.

A single-valued holomorphic function on $\Omega'$ algebraic over
$\mathbf C(t)$ belongs to $\mathbf C(t)$. To see this, normalize
the base in its irreducible function-field extension and remove
the remaining branch and pole locus. This gives a connected
finite \'etale cover. Its complex analytic space is connected:
an irreducible complex variety remains analytically connected
after deleting a proper algebraic subset. A globally defined
branch is a section of this covering and has open-and-closed
image, so the connected cover has degree one. The same reasoning
applies to the antiholomorphic coefficients after conjugation.

Finally a rational function in $\mathbf C(t)$ algebraic over
$\overline{\mathbf Q}(t)$ belongs to
$\overline{\mathbf Q}(t)$. One concrete descent argument is to
bound numerator and denominator degrees and choose sufficiently
many algebraic parameter points avoiding their poles and the
degenerate specializations of the algebraic equation. The values
are algebraic. Rational interpolation, normalized by one nonzero
coefficient, then recovers the function by a finite linear system
with algebraic entries. Normalize one denominator coefficient
known to be nonzero. A solution exists over $\mathbf C$, hence
also over $\overline{\mathbf Q}$ by Gaussian elimination.
Sufficiently many interpolation points determine the rational
function uniquely (cross-multiply and interpolate a polynomial
of twice the degree bound), so this algebraic solution represents
the original function. Consequently
\[
 H\in\overline{\mathbf Q}(t)[Z],\qquad
 A\in\overline{\mathbf Q}(\bar t)[Z].
\]

\paragraph{Equivariance eliminates every nonconstant anti factor.}
Choose a finite number field $M$ stable under the fixed conjugation
$c$, containing every coefficient and defining equation in the
argument. For $\sigma\in G_{EM}$, equivariance at algebraic
parameters and cancellation of the common monic factor give
\[
 A(c\sigma t,Z)=A(\sigma c t,Z).
\]
This holds on a nonempty Zariski open set: all required domains
are finite intersections of opens, and their transforms under
the field automorphisms are again nonempty Zariski opens. For
$\tau=c\sigma c\sigma^{-1}$ and $u=\sigma ct$, it becomes
$A(\tau u,Z)=A(u,Z)$. Both $\sigma$ and $c\sigma c$ fix $M$,
so $\tau$ fixes the coefficients of $A$. If one coefficient
$a(u)$ is nonconstant, the equality says
$\tau(a(u))=a(u)$ on that open set. Its image contains a cofinite
subset of $\overline{\mathbf Q}$: restrict to an algebraic line
on which the rational function is nonconstant and the open set
is nonempty, and use that a nonconstant rational map of the
projective line is surjective, with only finitely many deleted
source points. An automorphism fixing a cofinite subset fixes
every element, since every element is a difference of two
members of that subset. Hence $\tau=1$ for every
$\sigma\in G_{EM}$, contradicting B.B8.1(3).

All coefficients of $A$ are therefore constant. Its roots are
then constant, whereas by construction $A$ has no constant root
branch. Thus $A=1$ and $P=H$ is the required monic polynomial
over $\overline{\mathbf Q}(t)$.

\paragraph{Check outcome and limits.}
This targeted second reading verifies the mathematical argument
of B.B8.3, with the explicit continuation justification above.
It does not independently re-prove the Carlson--Haine--Wolf
pinned-point-map theorem, Baker--Brumer, or the henselian
valuation theorem used in B.B8.1. Those are separate source
inputs. It does not certify B.B8.4, which is being reviewed by
the parent researcher. The check is by the same advertised model
family and is not represented as cross-model validation.
